---

### File 2: `paper.tex`

```latex
\documentclass[11pt,a4paper]{article}

\usepackage[utf8]{inputenc}
\usepackage{amsmath,amssymb,amsthm}
\usepackage{geometry}
\usepackage{hyperref}
\usepackage{booktabs}

\geometry{margin=1in}

\theoremstyle{plain}
\newtheorem{theorem}{Theorem}
\newtheorem{lemma}{Lemma}
\newtheorem{corollary}{Corollary}

\theoremstyle{definition}
\newtheorem{definition}{Definition}
\newtheorem{remark}{Remark}

\title{\textbf{On the Bounded Growth, Streak Invariants, and Drift Thresholds of Generalized Collatz Mappings}}
\author{\textbf{[Jimmy Nakatsu]} \\ \small Independent Researcher}
\date{August 2026}

\begin{document}

\maketitle

\begin{abstract}
We examine a generalized parameter family of discrete dynamical systems $T: \mathbb{Z}^+ \to \mathbb{Z}^+$ extending the classical Collatz ($3x+1$) mapping over an arbitrary integer base $m \ge 2$, multiplier parameter $k \ge 1$, and exact residue offset $(m - r)$. We establish that while individual orbits may exhibit divergent behavior under expansive regimes, unbroken sequences of non-divisible iterations are strictly bounded by positional base-$m$ digit constraints. We formulate three foundational lemmas bounding single-step growth, digit capacity, and modular carry exhaustion, proving that an unbroken expansion streak for any $L$-digit seed satisfies $M \le L$ with a peak ceiling of $(m+1)^L - 1$. Furthermore, by evaluating the expected $m$-adic valuation of the affine branch, we prove that the global expected logarithmic drift is contractive if and only if $m + k < m^{\frac{m}{m-1}}$, demonstrating that $k=1$ is the unique contractive integer mapping in the binary ($m=2$) domain.
\end{abstract}

\section{Introduction}
The classical Collatz conjecture concerns the iterative behavior of the piecewise affine map on positive integers:
\begin{equation}
T_{\text{std}}(x) = \begin{cases} 
\frac{x}{2}, & x \equiv 0 \pmod 2 \\[4pt]
\frac{3x+1}{2}, & x \equiv 1 \pmod 2
\end{cases}
\end{equation}
While computational verifications have validated convergence for all $x_0 \le 2.95 \times 10^{20}$, formal proofs of global stability remain elusive. In this paper, we generalize the system to an arbitrary integer base $m \ge 2$ and investigate the interaction between modular residue consumption, finite streak ceilings, and global geometric drift.

\section{The Generalized Mapping Framework}

\begin{definition}[Generalized Collatz Map]
Let $m \in \mathbb{Z}_{\ge 2}$ and $k \in \mathbb{Z}^+$ such that $(k-1) \equiv 0 \pmod m$. Define $T: \mathbb{Z}^+ \to \mathbb{Z}^+$ by:
\begin{equation}
T(x) = \begin{cases} 
\dfrac{x}{m}, & x \equiv 0 \pmod m \\[8pt]
\dfrac{(m + k)x + (m - (x \bmod m))}{m}, & x \not\equiv 0 \pmod m
\end{cases}
\end{equation}
\end{definition}

The adjustment term $(m - r)$, where $r = x \bmod m \in \{1, 2, \dots, m-1\}$, ensures that the non-zero branch maps integers directly to integers for all valid remainder classes. Setting $m=2, k=1$ recovers the standard binary Collatz formulation.

\section{Finite Streak Invariants and Peak Ceilings}

\begin{lemma}[Single-Step Growth Envelope]
For any $x \in \mathbb{Z}^+$ with $x \not\equiv 0 \pmod m$:
\begin{equation}
T(x) \le \left(\frac{m+k}{m}\right)x + \frac{m-1}{m}
\end{equation}
Equality holds strictly if and only if $x \equiv 1 \pmod m$.
\end{lemma}

\begin{proof}
Let $r = x \bmod m \in \{1, \dots, m-1\}$. Expanding the definition:
\begin{equation}
T(x) = \left(\frac{m+k}{m}\right)x + 1 - \frac{r}{m}
\end{equation}
Since $r \ge 1$, we have $-\frac{r}{m} \le -\frac{1}{m}$, which yields $1 - \frac{r}{m} \le \frac{m-1}{m}$. The maximum is achieved uniquely at $r = 1$.
\end{proof}

\begin{lemma}[Base-$m$ Positional Representation Bound]
Any positive integer $x_0$ with $L$ digits in base $m$ satisfies:
\begin{equation}
m^{L-1} \le x_0 < m^L \implies x_0 \le m^L - 1
\end{equation}
\end{lemma}

\begin{lemma}[Modular Streak Termination]
Let $M$ denote the number of consecutive non-zero steps ($x_i \not\equiv 0 \pmod m$) starting from $x_0$. An unbroken streak of length $M$ requires:
\begin{equation}
x_0 \equiv m^M - 1 \pmod{m^M}
\end{equation}
Consequently, for any $x_0 < m^L$, the maximum unbroken streak length is bounded by $M \le L$.
\end{lemma}

\begin{proof}
Each non-zero expansion step consumes one trailing maximal digit $(m-1)$ in the base-$m$ representation through carry propagation. Since an integer $x_0 < m^L$ possess at most $L$ base-$m$ digits, the sequence of maximal carries must terminate in at most $L$ iterations, forcing a division step.
\end{proof}

\begin{theorem}[Maximum Peak Height for $L$-Digit Trajectories]
Let $x_0 \in \mathbb{Z}^+$ have $L$ digits in base $m$, and let $k=1$. The maximum value attainable after an unbroken streak of $M$ non-zero steps satisfies:
\begin{equation}
\max_{x_0 < m^L} x_M \le (m+1)^L - 1
\end{equation}
\end{theorem}

\begin{proof}
Applying Lemma 1 inductively over $M$ consecutive steps generates the recurrence:
\begin{equation}
x_M \le \left(\frac{m+1}{m}\right)^M (x_0 + m - 1) - (m - 1)
\end{equation}
Substituting the maximal initial seed $x_0 = m^L - 1$ from Lemma 2 and the streak bound $M \le L$ from Lemma 3 yields:
\begin{equation}
x_M \le \left(\frac{m+1}{m}\right)^L (m^L) - (m-1) = (m+1)^L - m + 1 \le (m+1)^L - 1
\end{equation}
\end{proof}

\section{Global Stopping Time Threshold}

\begin{theorem}[Critical Drift Threshold]
Let $\widetilde{T}(x) = \frac{(m+k)x + (m - r)}{m^{1 + d(x)}}$ denote the accelerated map, where $d(x) = \nu_m((m+k)x + (m-r))$ is the $m$-adic valuation. Under uniform modular distribution, the expected logarithmic drift $\mathbb{E}[\Delta \ln x]$ is negative if and only if:
\begin{equation}
m + k < m^{\frac{m}{m-1}}
\end{equation}
\end{theorem}

\begin{proof}
The probability of obtaining exactly $j$ additional divisions by $m$ is given by $\mathbb{P}(d = j) = \frac{m-1}{m^{j+1}}$ for $j \ge 0$. The expected total division exponent is:
\begin{equation}
\mathbb{E}[1 + d] = 1 + \sum_{j=1}^{\infty} \frac{1}{m^j} = 1 + \frac{1}{m-1} = \frac{m}{m-1}
\end{equation}
The expected logarithmic drift per cycle is:
\begin{equation}
\mathbb{E}[\Delta \ln x] = \ln(m+k) - \mathbb{E}[1+d]\ln(m) = \ln(m+k) - \left(\frac{m}{m-1}\right)\ln(m)
\end{equation}
Enforcing strict contractive drift $\mathbb{E}[\Delta \ln x] < 0$ yields $\ln(m+k) < \ln\left(m^{\frac{m}{m-1}}\right)$, from which exponentiation gives $m + k < m^{\frac{m}{m-1}}$.
\end{proof}

\begin{corollary}[Uniqueness of Binary Contraction]
For $m = 2$, the inequality reduces to $2 + k < 2^2 = 4 \implies k < 2$. The only positive integer solution is $k = 1$ ($3x+1$).
\end{corollary}

\begin{corollary}[Asymptotic Margin for Large Bases]
As $m \to \infty$, $m^{\frac{1}{m-1}} = 1 + \frac{\ln m}{m} + O\left(\frac{\ln^2 m}{m^2}\right)$. Thus, the allowable multiplier parameter satisfies:
\begin{equation}
k < \ln(m)
\end{equation}
\end{corollary}

\section{Conclusion}
We have shown that while localized growth streaks in generalized Collatz mappings are strictly bounded by positional carry invariants ($M \le L$), global contractive drift requires parameter adherence to $m + k < m^{\frac{m}{m-1}}$. This classification formalizes why standard $3x+1$ occupies a unique stable position in base-2 arithmetic dynamics.

\end{document}
